\documentclass[11pt]{article}

\usepackage[margin=1in]{geometry}
\usepackage{amsmath,amssymb}
\usepackage{booktabs}
\usepackage{hyperref}

\begin{document}

\subsection{Best PF-identity GNN architectures per dataset}
\label{sec:best-arch-per-dataset}

Across all three full-voltage datasets---Original, Derived (injection), and Load-type---the final selected models use the same \emph{PF-identity GNN} backbone. The datasets differ only in the node input features and in a small set of architectural hyperparameters, namely the node and edge embedding dimensions $(n_{\text{emb}}, e_{\text{emb}})$, hidden width $h$, depth $L$, and whether explicit phase encoding is used.

For a given dataset, let $x_i \in \mathbb{R}^{d_{\text{node}}}$ denote the feature vector at bus-phase node $i$. Each directed edge $(j \to i)$ carries physical line parameters
\[
e_{ji} = [R_{ji}, X_{ji}] \in \mathbb{R}^2
\]
together with an integer edge index. The model predicts the per-unit voltage magnitude $v_i$ at every node.

\subsubsection*{Common PF-identity GNN backbone}

The shared architecture is defined as follows.

\paragraph{1. Node and edge embeddings.}
Each node $i$ has a learned embedding $z_i \in \mathbb{R}^{n_{\text{emb}}}$ and each edge $(j,i)$ has a learned embedding $r_{ji} \in \mathbb{R}^{e_{\text{emb}}}$:
\[
z_i = \mathrm{Emb}_{\mathrm{node}}(i), 
\qquad
r_{ji} = \mathrm{Emb}_{\mathrm{edge}}(\mathrm{edge\_id}(j,i)).
\]

\paragraph{2. Input encoding.}
The initial hidden state is obtained by concatenating the node features with the node embedding and passing the result through an MLP encoder:
\[
h_i^{(0)} = \phi_0([x_i, z_i]).
\]
Here,
\[
\phi_0 : \mathbb{R}^{d_{\text{node}} + n_{\text{emb}}} \to \mathbb{R}^{h}
\]
is a 2-hidden-layer MLP of the form
\[
\mathrm{Linear}(d_{\text{node}}+n_{\text{emb}},h)
\to \mathrm{ReLU}
\to \mathrm{Linear}(h,h)
\to \mathrm{ReLU}
\to \mathrm{Linear}(h,h).
\]

\paragraph{3. Message passing.}
For each layer $\ell = 0,\dots,L-1$, a message is computed on every directed edge $(j \to i)$:
\[
m_{ji}^{(\ell)} =
\psi^{(\ell)}\!\left([h_j^{(\ell)}, e_{ji}, r_{ji}]\right),
\]
where
\[
\psi^{(\ell)} : \mathbb{R}^{h+2+e_{\text{emb}}} \to \mathbb{R}^{h}
\]
is a 2-hidden-layer MLP with hidden width $h$. Incoming messages are summed,
\[
m_i^{(\ell)} = \sum_{j \in \mathcal{N}(i)} m_{ji}^{(\ell)},
\]
and the node state is updated as
\[
h_i^{(\ell+1)} =
\mathrm{Upd}^{(\ell)}\!\left([h_i^{(\ell)}, m_i^{(\ell)}, z_i]\right),
\]
where
\[
\mathrm{Upd}^{(\ell)} : \mathbb{R}^{2h+n_{\text{emb}}} \to \mathbb{R}^{h}
\]
is again a 2-hidden-layer MLP with hidden width $h$.

\paragraph{4. Readout.}
After $L$ message-passing layers, the per-node voltage magnitude is predicted by
\[
\hat v_i = \rho(h_i^{(L)}),
\qquad
\rho : \mathbb{R}^{h} \to \mathbb{R},
\]
where $\rho$ is a final 2-hidden-layer MLP.

\paragraph{5. Training objective.}
All selected models are trained using mean-squared error on voltage magnitude:
\[
\mathcal{L} = \frac{1}{N}\sum_{i=1}^{N}(\hat v_i - v_i)^2.
\]
Training uses Adam, early stopping based on validation RMSE, and a fixed split of snapshots into training and test sets.

\subsubsection*{Architecture search summary}

The same PF-identity backbone was used throughout the GNN2 and GNN3 studies, while varying depth, width, embedding size, normalization, and phase encoding.

For the Original and Derived datasets, the search included shallow and deep variants ranging from 2 to 8 message-passing layers, with hidden widths spanning approximately $32$ to $192$, node/edge embedding sizes such as $(8,4)$ and $(16,8)$, and selected normalized or phase-aware variants. In both datasets, deeper models beyond 4 layers did not improve performance, and the best models were either moderately deep and wide or shallow and wide.

For the Load-type dataset, the architecture search was broader. In addition to width/depth sweeps, it included normalization, phase one-hot encodings, and phase-subgraph variants. The final best model on this dataset emerged from the refine stage and showed that explicit phase encoding materially improves accuracy.

\subsubsection*{Dataset-specific best models}

\paragraph{Original dataset (baseline PF outputs).}
For the Original dataset, the node features are
\[
x_i =
\big[
p_{\text{load},i}^{(\mathrm{kW})},
\;
q_{\text{load},i}^{(\mathrm{kvar})},
\;
p_{\text{pv},i}^{(\mathrm{kW})},
\;
q_{\text{pv},i}^{(\mathrm{kvar})}
\big]
\in \mathbb{R}^{4}.
\]
Using the current GNN2 architecture search (on one third of the snapshots with early stopping), the best-performing model is the \emph{light\_wide\_emb\_depth3} configuration, with
\[
n_{\text{emb}} = 12,\qquad
e_{\text{emb}} = 6,\qquad
h = 96,\qquad
L = 3.
\]
This model uses no normalization and no explicit phase encoding. It achieves
\[
\mathrm{MAE} \approx 0.00583,
\qquad
\mathrm{RMSE} \approx 0.00781\;\text{pu}.
\]
Among the Original-dataset candidates in this search, this moderately wide 3-layer model outperformed both shallower and deeper alternatives on the held-out test set.

\paragraph{Derived (injection) dataset.}
For the Derived dataset, the node features encode net complex power injection:
\[
x_i =
\big[
p_{\text{inj},i}^{(\mathrm{kW})},
\;
q_{\text{inj},i}^{(\mathrm{kvar})}
\big]
\in \mathbb{R}^{2}.
\]
Under the same GNN2 search protocol, the best-performing model is the \emph{wide\_shallow\_h128\_depth3} configuration, with
\[
n_{\text{emb}} = 8,\qquad
e_{\text{emb}} = 4,\qquad
h = 128,\qquad
L = 3.
\]
This model also uses no normalization and no explicit phase encoding. It achieves
\[
\mathrm{MAE} \approx 0.00567,
\qquad
\mathrm{RMSE} \approx 0.00771\;\text{pu}.
\]
Compared to deeper 4-layer alternatives and wider variants, this shallow but reasonably wide configuration gives the lowest error on the injection dataset.

\paragraph{Load-type dataset.}
For the Load-type dataset, the base node features now consist of 14 scalar descriptors:
\[
x_i =
\big[
\text{electrical distance},\;
m1\_p,\; m1\_q,\;
m2\_p,\; m2\_q,\;
m4\_p,\; m4\_q,\;
m5\_p,\; m5\_q,\;
q_{\text{cap}},\;
p_{\text{pv}},\;
q_{\text{pv}},\;
p_{\text{sys\_balance}},\;
q_{\text{sys\_balance}}
\big]
\in \mathbb{R}^{14}.
\]
Some candidate models augment these features with a one-hot phase indicator
$(\mathbf{1}\{\phi_i=A\},\mathbf{1}\{\phi_i=B\},\mathbf{1}\{\phi_i=C\})$, yielding a 17-dimensional input. In the current GNN2 search, however, the best-performing configuration is the simpler
\emph{light\_emb\_h96\_depth2}, with
\[
n_{\text{emb}} = 16,\qquad
e_{\text{emb}} = 8,\qquad
h = 96,\qquad
L = 2.
\]
This model uses no normalization and does \emph{not} include explicit phase one-hot encoding. It achieves
\[
\mathrm{MAE} \approx 0.00158,
\qquad
\mathrm{RMSE} \approx 0.00221\;\text{pu}.
\]
This is the best result among the full-voltage Load-type experiments in the current search, indicating that a relatively shallow architecture with accurate local features (including $q_{\text{pv}}$) is sufficient once the dataset and feature definitions are aligned with the present codebase.

\subsubsection*{Final selected architectures}

Table below summarize the final selected architecture for each non--$\Delta V$ full-voltage dataset.

\begin{table}[htbp]
  \centering
  \footnotesize
  \caption{Best-performing PF-identity GNN architectures for all full-voltage datasets.}
  \label{tab:best-gnn-arch}
  \begin{tabular}{llllllllll}
    \toprule
    Dataset & Config & $n_{\text{emb}}$ & $e_{\text{emb}}$ & $h$ & $L$ & Phase enc. & MAE / RMSE (pu) \\
    \midrule
    Original 
        & light\_wide\_emb\_depth3
        & 12 & 6 & 96 & 3 & none  
        & $0.00583 / 0.00781$ \\
    Derived  
        & wide\_shallow\_h128\_depth3
        & 8 & 4 & 128 & 3 & none 
        & $0.00567 / 0.00771$ \\
    Load-type (no global) 
        & light\_emb\_h96\_phase\_onehot\_depth3\_h112\_no\_global
        & 16 & 8 & 112 & 3 & one-hot 
        & $0.00537 / 0.00733$ \\
    Load-type (no embeddings) 
        & light\_h96\_depth2\_noemb
        & 0 & 0 & 96 & 2 & none
        & $0.00683 / 0.00810$ \\
    Load-type 
        & light\_emb\_h96\_depth2
        & 16 & 8 & 96 & 2 & none
        & $0.00158 / 0.00221$ \\
    \bottomrule
  \end{tabular}
\end{table}

\end{document}